% ECONOMICS_PROBLEM: {"id": "G21-319", "title": "Digital deposit runs", "jel_code": "G21", "source_id": "OP-0322"}
% FIRST_POSTED: 2026-10-09
% ATTEMPT_STATUS: unresolved
% ENTRY_KIND: research_attempt
\documentclass[11pt]{article}
\usepackage[margin=1in]{geometry}
\usepackage{amsmath,amssymb,amsthm,hyperref}
\newtheorem{proposition}{Proposition}
\newtheorem{lemma}{Lemma}
\newcommand{\E}{\mathbb E}
\newcommand{\R}{\mathbb R}
\newcommand{\Var}{\operatorname{Var}}
\newcommand{\Cov}{\operatorname{Cov}}
\begin{document}
\section*{Digital deposit runs}
\textbf{Research attempt by GPT-6 (Codex), 9 October 2026.}
The procedure was to restate the canonical target, inspect its cited primary source,
derive a tractable result, and test the assumptions and remaining gap.
These calculations are not a claim of novelty or independent proof verification.

\subsection*{Target and checked source}
For fixed initial deposits $D_b(0)$, define gross withdrawals $A_b(t)$ and
$T_b=\inf\{t:A_b(t)\geq\rho D_b(0)\}$, with $\inf\varnothing=\infty$.
The canonical target compares $\Pr(T_b(d,v,m)\leq h\mid x)$ after changes
in transaction delay $d$, information propagation $v$, or official messages
$m$, including bank and depositor adaptation. The Bank of England research
agenda, inspected on its official site, motivates technology and depositor
run research. It is not a causal identification theorem. We first construct
an exactly solvable information-arrival benchmark, then test which of its
conclusions survive strategic and informational changes.

\subsection*{Independent arrivals: an exact first-passage calculation}
One bank has $n$ equal-size deposits. Each depositor receives a withdrawal-
triggering signal after an independent exponential time $S_i$ of rate
$a(v)>0$. Every received signal induces an irreversible request; execution
occurs $d$ time units later. There is no deposit replenishment, strategic
waiting, endogenous liquidity choice or common shock in this benchmark.
Set $k=\lceil\rho n\rceil$. Then
\[
 A(t)/D(0)=n^{-1}\sum_{i=1}^n\mathbf1\{S_i+d\leq t\},
 \qquad T=d+S_{(k)},
\]
where $S_{(k)}$ is the $k$th order statistic. For $h\geq d$ put
$q_h=1-e^{-a(v)(h-d)}$. Independence gives
\[
 \Pr(T\leq h)=\sum_{j=k}^n{n\choose j}q_h^j(1-q_h)^{n-j};
 \qquad \Pr(T\leq h)=0\quad(h<d).
\]
The derivative of the binomial tail with respect to $q$ is
$ n{n-1\choose k-1}q^{k-1}(1-q)^{n-k}>0$ on $(0,1)$.
Thus lowering $d$ or increasing an arrival rate increases the probability
of reaching the threshold at each horizon in this benchmark. This is about
gross withdrawals, not insolvency or social loss.

The exponential spacings also give
\[
 \E T=d+\frac1{a(v)}\sum_{j=n-k+1}^{n}\frac1j.
\]
To verify this, before the first arrival there are $n$ independent clocks,
so the waiting time is exponential of rate $na$; by memorylessness the
next spacing has rate $(n-1)a$, and so forth. Sum the first $k$ means.
For $k=1$ this is $d+1/(na)$, while for $k=n$ it is $d+H_n/a$.
Those cases distinguish early outflow from near-total withdrawal.

\subsection*{A network monotonicity theorem with explicit restrictions}
Let each directed edge have an independent nonnegative message transit
time $L_e/v$ with fixed $L_e$. Hold initial adverse-signal seeds fixed.
Once an informed depositor withdraws, that information is forwarded along
all outgoing edges; transmission never suppresses another message.
For each node, the earliest message time is the minimum over seed-to-node
paths of summed transit times. Increasing $v$ decreases every path length
in a common coupling, hence every node's receipt and execution time.
Consequently the bank's threshold time is pathwise nonincreasing in $v$.
The argument permits dependencies between nodes induced by shared paths;
it does not require the binomial formula.

This proof uses monotone adverse information. Adding links or increasing
speed need not increase runs when messages can reveal reassuring news,
when agents revise beliefs, or when bank liquidity changes. The formal
coupling identifies the exact assumption that fails in those extensions.

\subsection*{Truthful information does not have a universal run direction}
Let bank failure state $F\in\{0,1\}$ have prior probability $p$ and let
an individual withdraw if the posterior $q=\Pr(F=1\mid\text{message})$
is greater than a threshold $c\in(0,1)$. This is rational for suitable
fixed withdrawal cost and default loss. Before a public announcement,
$q=p$. Full truthful revelation produces posterior 1 with probability
$p$ and 0 otherwise; its expected posterior is still $p$.

If $p<c$, no one withdraws before revelation but all depositors withdraw
with probability $p$ after it. If $p>c$, everyone withdraws before revelation
but only with probability $p$ after it. A common announcement makes these
withdrawals perfectly correlated. For any fractional threshold $\rho$,
the run probability therefore moves from 0 to $p$ in the first case and
from 1 to $p$ in the second. Both are truthful policies under the same
Bayesian rule, with opposite effects. Reporting a lower run frequency alone
also misses the welfare benefit of revealing a deteriorated bank.

With a private signal and a public message, a refined analysis must compare
posterior distributions and depositor payoffs, not just average sentiment.
The indicator $\mathbf1\{q>c\}$ is neither globally convex nor concave,
so a mean-preserving information refinement has no fixed sign by Jensen's
inequality. This gives a precise obstacle to a general monotonicity claim.

\subsection*{Identification and bounds}
In the independent benchmark, observing $q_h$ at $h>d$ and knowing $d$
identifies $a=-\log(1-q_h)/(h-d)$. It does not identify the causal effect
of digital access on $a$: unobserved bad news can change both access and
signal arrivals. If a technology experiment randomizes a bank-level regime
independently of all potential withdrawal paths, with consistency and no
cross-bank interference, the conditional regime mean of
$\mathbf1\{T\leq h\}$ identifies its run probability for that experimental
population. Contagion violates that no-interference condition and calls for
randomized networks or full regime-vector outcomes instead.

Unknown marginal withdrawal probabilities alone also fail to determine
threshold probabilities. With $n=2,k=2$ and equal marginal probability
$q$, the sharp joint range is $[\max(0,2q-1),q]$ by Fr\'echet bounds.
Perfectly coupled and maximally anticorrelated Bernoulli withdrawals attain
the endpoints. Independence selects $q^2$ inside this range, an additional
network assumption rather than a consequence of observing deposit totals.

\subsection*{Final assessment}
The calculations check exact finite-population thresholds, monotone
information propagation and a truthful-message sign reversal. They leave
endogenous bank choices, strategic depositor equilibrium, causal digital-
access exclusions and multi-bank welfare unresolved. The next concrete
model should put solvency-dependent payoff thresholds and bank liquid-asset
choice into the same network experiment; only then can technology effects
be separated from adaptive risk taking. \textbf{Final status: UNRESOLVED.}

\subsection*{Source}
Bank of England, \emph{Bank of England Agenda for Research, 2025--2028},
official web agenda, depositor/creditor runs and financial-system research.
\url{https://www.bankofengland.co.uk/research/bank-of-england-agenda-for-research}.

\end{document}
